% ==============================================================================
% ================================== AUFGABE 4 ================================
% ==============================================================================

\chapter{Aufgabe 4: Handlungsalternativen bei Scheitern des Zweitgesprächs} 
\label{sec:aufgabe4}

\section{Einordnung: Erreichen der Grenzen der Selbsthilfe}
Sollte auch das zweite persönliche Klärungsgespräch zwischen dem Dienstleister und der Kundin ohne Konsens oder konstruktives Ergebnis verlaufen, ist die Ebene der bilateralen Selbsthilfe endgültig überschritten. Die Beteiligten befinden sich in einem Zustand der Verstrickung, in dem sachliche Differenzen und persönliche Beziehungsstörungen einander wechselseitig verstärken \cite[S.~42--43]{5}. 

In dieser Eskalationsphase greift der Grundsatz, dass die Beteiligten den Konflikt nicht mehr eigenständig steuern können, sondern vom Konflikt gesteuert werden („Der Konflikt hat uns“) \cite[S.~42]{5}. Zur Vermeidung einer weiteren Eskalation bis hin zum gerichtlichen Parteienstreit oder der totalen Beziehungszerstörung müssen systematische Drittinterventionen und professionelle Verfahrensmodelle herangezogen werden \cite[S.~56]{5}.

Kiefer ordnet den Konflikt nach dem neunstufigen Eskalationsmodell von Glasl auf Stufe 5 (Gesichtsverlust) bis Stufe 6 (Drohstrategien) ein: Die Kundin hat den Dienstleister im Gespräch massiv beschimpft, sein Gesicht verletzt und ihm ein Ultimatum („bis heute Abend“) gesetzt (Glasl, zit. n. \cite[S.~34--38]{5}). Auf dieser Eskalationsstufe ist eine bilaterale Selbsthilfe nicht mehr möglich; die moderne Mediation (Stufen 1–7) bietet sich als primäres Verfahren an \cite[S.~70--72]{5}.

\section{Systematische Handlungsalternativen}

\subsection{Wirtschaftsmediation (Einbindung eines neutralen Mediators)}
Bevor ein externer Mediator hinzugezogen wird, kann ein neutraler Dritter aus dem eigenen Umfeld (z. B. ein gemeinsamer Projektleiter oder Key Account Manager) als Nachbarschaftshelfer vermitteln \cite[S.~53--54]{5}.
Die primäre Handlungsalternative bei festgefahrenen Verhandlungen stellt die Wirtschaftsmediation dar. Bei diesem Verfahren führt ein unabhängiger, neutraler Dritter (Mediator) die Konfliktparteien durch einen strukturierten Verhandlungsprozess nach den Vorgaben des Mediationsgesetzes.

\begin{itemize}
    \item \textbf{Eignung und Nutzen:} Die Mediation setzt auf das Konsensprinzip und strebt eine zukunftsorientierte Vereinbarung an, von der beide Seiten profitieren \cite[S.~70--71]{5}.
    \item \textbf{Anwendung im konkreten Fall:} Ein externer Wirtschaftsmediator schafft einen geschützten Rahmen, in dem die tiefgreifenden Ursachen (Vorstandsdruck auf Kundenseite sowie fehlende Zuarbeiten) ohne gegenseitige Schuldzuweisungen offengelegt werden können. Durch die Allparteilichkeit des Mediators wird das Machtungleichgewicht abgebaut und die Basis für ein Nachbesserungskonzept erarbeitet \cite[S.~73--74]{5}.
\end{itemize}

\subsection{Technisches Gutachter- und Schiedsverfahren (Arbitration)}
Können sich die Parteien nicht auf ein freiwilliges Konsensverfahren wie die Mediation einigen oder bleibt die Klärung komplexer Fachfragen strittig, bietet sich die Einschaltung eines unabhängigen IT-Sachverständigen im Rahmen eines Schiedsgutachtenverfahrens an \cite[S.~63]{5}.

\begin{itemize}
    \item \textbf{Eignung und Nutzen:} Im Unterschied zur Mediation erarbeitet der Schiedsgutachter keine eigenständige Beziehungslösung, sondern trifft eine verbindliche Feststellung bezüglich des technischen Ist-Zustands und der Vertragsgegebenheiten.
    \item \textbf{Anwendung im konkreten Fall:} Der Gutachter prüft neutral, ob die bisherigen Automatisierungsmodule den Anforderungen des vereinbarten Lastenhefts entsprechen und inwiefern das Fehlen der Testdaten die Abweichungen verursacht hat. Die Parteien unterwerfen sich vorab vertraglich dem fachlichen Urteil des Gutachters.
\end{itemize}

\subsection{Hierarchischer Machteingriff (Eskalation an das Management)}
Vor der Eskalation zum Kundenmanagement sollte der Dienstleister seinen eigenen Vorgesetzten einbinden, um die Strategie abzustimmen und abzusichern (Nachbarschaftshilfe im erweiterten Sinne, \cite[S.~53--54]{5}).
Verweigert die Kundin die Mitwirkung an außergerichtlichen Schlichtungsverfahren, bleibt als innerorganisatorische bzw. vertragliche Alternative die methodische Eskalation an das übergeordnete Management (z.~B. Lenkungsausschuss oder Vorstand) \cite[S.~62--63]{5}.

\begin{itemize}
    \item \textbf{Eignung und Nutzen:} Ein Machteingriff erfolgt durch eine übergeordnete Instanz, die die Befugnis besitzt, verbindliche Weisungen zur Konfliktbeilegung durchzusetzen.
    \item \textbf{Anwendung im konkreten Fall:} Der Dienstleister wendet sich sachlich und unter Vorlage der Dokumentation (E-Mail-Protokolle, Lastenheft) an das Projekt-Steuerungsgremium des Kunden. Das Gremium kann der Kundin verbindliche Ressourcen und Fristen zur Bereitstellung der Testdaten zuweisen oder das Projekt offiziell umpriorisieren \cite[S.~64]{5}.
\end{itemize}

\subsection{Geordnete Vertragsanpassung oder Projektterminierung (Exit-Strategie)}
Sollten sämtliche Schlichtungs- und Eskalationsversuche scheitern und die Kooperationsbasis unheilbar zerstört sein, besteht die letzte Handlungsoption in der geordneten Vertragsauflösung (Exit-Strategie) \cite[S.~62]{5}.

\begin{itemize}
    \item \textbf{Eignung und Nutzen:} Ziel ist die geordnete Schadensabwicklung und der Schutz des Dienstleisters vor unkalkulierbaren Haftungs- und Folgekosten.
    \item \textbf{Anwendung im konkreten Fall:} Der Dienstleister übergibt den bisher entwickelten Quellcode im aktuellen Ist-Zustand gegen Teilvergütung der erbrachten Leistungen. Die Zusammenarbeit wird einvernehmlich beendet, um langwierige juristische Auseinandersetzungen zu vermeiden.
\end{itemize}

\section{Bewertung und Handlungsempfehlung}
Aus Sicht des Konfliktmanagements ist die \textbf{Wirtschaftsmediation} als bevorzugte Option zu empfehlen, da sie als einziges Verfahren die Chance bietet, sowohl das Automatisierungsprojekt erfolgreich fertigzustellen als auch die Kundenbeziehung langfristig zu erhalten. Erst wenn die Freiwilligkeit zur Mediation auf Kundenseite fehlt, sollten gestufte Maßnahmen über das Schiedsgutachten oder die Management-Eskalation eingeleitet werden.

% ==============================================================================
% ================================== AUFGABE 4 ================================
% ==============================================================================